\subsection{增益来源}

本文一项最稳定的结果是，在这首 29k 字符的叙事诗上，答案正确率主要取决于原始诗节的召回，而非图结构或 LLM 画像。C1（纯 chunk 检索）已达 0.820，C2 仅高 0.025；去掉原文的 C5 降至 0.445，拒答率 47\%。C2 的 gold 证据召回率为 0.78，反而低于 C1 的 0.86，因为图画像占用了约 1,784 token 的上下文预算，原文仅剩约 1,064 token，而 C1 的 2,862 token 几乎全为原文。图结构确实提高了 citation recall（C2 0.61 对 C1 0.54），使更多可引用关系进入上下文，但这些关系未转化为答案正确率。

这与 LightRAG 原文的动机并不矛盾。原文面向百万 token 的技术或法律语料，chunk 检索容易碎片化；本文语料只有 28 个 chunk，top-8 已覆盖 86\% 的 gold 证据，图的一跳扩展（one-hop expansion）很难再提供边际信息。就此而言，在本语料规模上，图结构与纯 chunk 检索不可区分，并非无效。C7（打乱一跳结构）与 C2 的正确率几乎一致，可用同一机制解释。当匹配到的实体已覆盖全部相关节点时，邻居是否随机化对答案影响很小。

\subsection{结构与画像文本的归因}

H1 的四个关键对照（C2$-$C6、C2$-$C7、C6$-$C1、C2$-$C1）均不显著，仅 C2$-$C5 显著。这是一个清晰的负结果，没有证据表明 LLM 画像文本或图结构单独贡献了可检出的正确率增益。C6（空画像）与 C2 的差仅 0.005，C7（打乱结构）与 C2 的差仅 0.015，C6 相对 C1 也仅 0.020；C5 去掉原文后损失 0.400。LightRAG 的增益主要来自检索到原文，而非将原文改写为图。

该结论支持该博客的第一项批评，但需限定范围。它支持的是“在该语料上，图结构单独产生该量级增益仍是一个未被证实的假设”，而非“图结构在任何语料上都没有用”。本文未重复 UltraDomain 的百万 token 规模，不能外推；但至少说明，原文条件的胜率不能直接归因于图结构，须将原文、画像文本与结构三者分开。

\subsection{字面漂移成立与事实偏差未证实}

H2 的结果比该博客中的强版本判断更复杂。自动审计显示，LLM 画像的实体支持率为 0.952，高于抽取式画像的 0.883；关系支持率相当（0.871 对 0.877）；答案层的 citation precision、unsupported rate 与正确率均无显著差异。这与“LLM 改写一定引入更多事实错误”的预期不符。

字面一致性确实下降。LLM 画像的 ROUGE-L 为 0.020/0.026（实体/关系），抽取式画像为 0.030/0.055；字符相似度为 0.012/0.018 对 0.035/0.071。Hermaphroditus 的画像还提供了一个具体改写错误，即 “though not explicitly mentioned in the text”，而诗中明确写有 Hermaphroditus。这说明自动审计的高支持率不能替代逐条人工核查，30 条抽样可能漏掉低频但严重的错误。H2 的合理结论是字面漂移成立、事实偏差未获统计证实、个案风险存在，而非“改写一定更差”。

\subsection{标准图不弱与时序结构的定向收益}

H3a 不成立。C2 的时序题正确率（0.88）不低于其主题题（0.78），也不低于事实与多跳题，标准图并未在时序题上普遍失效。图相对无图有时序定向收益，C2 比 C1 高 0.200（$p=0.005$），说明图结构将分散在各诗节的时间线索关联起来。进一步加入事件节点与 \texttt{before} 边（C9）使时序题正确率达到 1.000，相对 C2 仅边缘显著（$p=0.077$，Holm $p=0.154$），相对只做排序的 C10（0.680）则显著（$p<0.001$）。关键在于时序结构，而非按序排列 chunk。

C9 的 1.000 需谨慎解读。时间线由人工整理，包含 55 个事件、54 条边，与时序 gold 的判断同源；它反映表示能力上限，并非自动抽取结果。C10 的失败（0.680）也说明，仅将 chunk 按诗节序重排并不能让模型理解事件顺序；时序图的价值在于显式事件与边，而非排序本身。

\subsection{评测与部署}

H4 得到支持。人工盲审与第一评判模型的二次加权 $\kappa$（quadratic weighted $\kappa$）为 0.797，严格二分类 $\kappa$ 为 0.535；评判模型平均 1.57，人工为 1.30，系统性偏宽松，在闭卷 C0 与空画像 C6 上最明显。若只看 LLM 评判模型的绝对分数，会高估方法质量；条件排序大体保留（Spearman 0.845），小样本、单评判模型的结论仍需 CI 与人工校准。本文所有主指标均带 bootstrap CI，并对 H1--H3 做 Holm 校正，较只报四维胜率更保守。

成本上，索引 28 次调用、32.5k token、调用墙钟合计 195.8 秒；C2 每 query 约 3.6k token、P50 4.5 秒、P95 6.4 秒；图条件每 query 两次 LLM 调用。原文的“检索阶段 <100 token、一次调用”只覆盖关键词抽取与匹配，不覆盖答案上下文，全生命周期成本要高一个量级。增量更新节省 50.9\% token，是 LightRAG 的真实优势；在这类小语料上，索引的绝对成本不高，节省的绝对量也有限。

\subsection{对实践的启示}

\begin{enumerate}
\item 在小语料上，先建立 chunk 检索基线；若 top-$k$ 已覆盖大部分 gold 证据，图结构的边际收益可能不可检出，不宜默认图结构更优。
\item 若使用图，应为原文保留固定预算。本文预留 1,200 token 原文是必要的；否则 LLM 画像会占用原文所需的上下文预算，C2 证据召回低于 C1 即为一例。
\item 画像优先用抽取式或可追溯的摘要。LLM 画像的平均支持率不一定低，字面漂移与个案错误仍需逐条审计；至少保留逐字来源。
\item 时序任务需要显式事件与边，而不是仅按诗节序排列 chunk；时序标注成本高，宜作为必要时的增强，而非默认组件。
\item 评测要报告 CI、人工校准与全生命周期成本；单评判模型的绝对分数会偏宽松，在拒答与部分正确的边界上尤其如此。
\end{enumerate}